\chapter{Two phases persist under perimeter blow-up}\label{sec:perimeter-two-phase}
A compactness limit can lose an entire phase unless both phases have quantitative volume. We derive the needed bounds from the already proved first-power Sobolev inequality. This expands \cite{Sim18}, Chapter 3, Theorem 4.4, printed p.~99, and avoids using the De Giorgi structure theorem in its own blow-up proof. Here $n\geq2$.

\begin{proposition}[The first-power inequality for compactly supported BV functions]\label{prop:bv-first-power-sobolev}
If $u\in BV(\mathbb R^n)$ has compact support, then
\[
 \|u\|_{L^{n/(n-1)}(\mathbb R^n)}\leq |Du|(\mathbb R^n).
\]
The constant is the convenient nonsharp constant in Theorem \ref{thm:sobolev-first-power}. This is the BV version of the Sobolev input in \cite{Sim18}, p.~99.
\end{proposition}
\begin{proof}
The smooth mollifications with a nonnegative kernel from Lemma \ref{lem:bv-mollification} have compact support, converge to $u$ in $L^1$ and have gradient integral at most $|Du|(\mathbb R^n)$. Apply Theorem \ref{thm:sobolev-first-power} to each mollification. Choose a sequence tending to $u$ almost everywhere and apply Fatou to the power $n/(n-1)$. Taking the appropriate root proves the assertion, including the conclusion that this higher-power norm is finite.
\end{proof}

\begin{corollary}[A cutoff Sobolev estimate for a phase]\label{cor:phase-cutoff-sobolev}
If $E$ has locally finite perimeter in $U$ and $0\leq\phi\in C_c^1(U)$, then
\[
 \left(\int_E\phi^{n/(n-1)}\,dx\right)^{(n-1)/n}
 \leq\int\phi\,d|D\chi_E|+\int_E|D\phi|\,dx.
\]
No bound on the total volume of $E$ is required. This is the cutoff inequality before equation (10) in \cite{Sim18}, p.~99.
\end{corollary}
\begin{proof}
By Proposition \ref{prop:bv-cutoff-product}, $\phi\chi_E$, extended by zero, is compactly supported and belongs to $BV(\mathbb R^n)$. Its variation is at most the sum on the right. Proposition \ref{prop:bv-first-power-sobolev} applies, and $|\phi\chi_E|^{n/(n-1)}=\phi^{n/(n-1)}\chi_E$. This proves the claim using only the distributional product rule.
\end{proof}

\begin{proposition}[A radial volume-perimeter inequality]\label{prop:phase-radial-sobolev}
If $B_R(x)\Subset U$ and $V(r)=|E\cap B_r(x)|$, then for almost every $0<r<R$,
\[
 V(r)^{(n-1)/n}\leq |D\chi_E|(B_r(x))+V'(r).
\]
The same inequality holds for $W(r)=|B_r(x)\setminus E|$ with the same perimeter measure. This is the radial cutoff passage on printed p.~99 of \cite{Sim18}.
\end{proposition}
\begin{proof}
Use the smooth radial functions $\phi_\delta$ constructed in Lemma \ref{lem:perimeter-radial-flux} in Corollary \ref{cor:phase-cutoff-sobolev}. Their values lie between zero and one and tend to the indicator of the closed ball. Its sphere is volume-null, so dominated convergence makes the left side tend to $V(r)^{(n-1)/n}$. At radii whose spheres have zero perimeter measure, the first term on the right tends to the open-ball perimeter. At Lebesgue points of $V'$, the second tends to $V'(r)$ by the weighted radial differentiation in that lemma. Both conditions hold almost everywhere. Apply the same argument to the complement for the final assertion.
\end{proof}

\begin{lemma}[A perimeter-support point sees positive volume of both phases]\label{lem:perimeter-support-two-phases}
If $x\in\operatorname{supp}|D\chi_E|$ and $B_r(x)\Subset U$, then
\[
 |E\cap B_r(x)|>0,\qquad |B_r(x)\setminus E|>0.
\]
This is a qualitative support assertion; it does not yet give a uniform density bound.
\end{lemma}
\begin{proof}
If the first volume were zero, $\chi_E=0$ almost everywhere on the ball. Its distributional derivative would vanish there by testing the defining integral, contradicting membership in the support. If the second volume were zero, use $\chi_E=1$ and the fact that the integral of the divergence of a compactly supported field is zero. In both cases the vector derivative and hence its total variation vanish on the ball, which is impossible.
\end{proof}

\begin{lemma}[Integrating a positive radial differential inequality]\label{lem:phase-volume-ode}
Suppose $V:[0,R)\to[0,\infty)$ is absolutely continuous on compact subintervals, $V(0)=0$, $V(r)>0$ for $r>0$, and
\[
 V(r)^{(n-1)/n}\leq C V'(r)\quad\text{for almost every }r\in(0,R),
\]
where $C>0$. Then $V(r)\geq(r/(nC))^n$ for every $0<r<R$.
\end{lemma}
\begin{proof}
The assumed inequality makes $V'\geq0$ almost everywhere, so $V$ is nondecreasing. On $[\delta,r]$ with $0<\delta<r<R$, it is bounded away from zero. The absolutely continuous chain rule therefore gives
\[
 \frac{d}{ds}V(s)^{1/n}=\frac1n V(s)^{-(n-1)/n}V'(s)\geq\frac{1}{nC}
\]
almost everywhere on that interval. Integrate, and let $\delta\downarrow0$. Absolute continuity at zero gives $V(\delta)\to0$; alternatively the nonnegative term can be discarded before taking this limit. Raising the resulting root inequality to the power $n$ proves the claim at every radius.
\end{proof}

\begin{theorem}[Two-sided phase volume at a reduced-boundary point]\label{thm:reduced-boundary-two-phase-volume}
If $x\in\partial^*E$, then for some $r_0>0$ and all $0<r<r_0$,
\[
 |E\cap B_r(x)|\geq c_n r^n,\qquad
 |B_r(x)\setminus E|\geq c_n r^n,
 \qquad c_n=(3n)^{-n}.
\]
The constant is not sharp. These are the two estimates (10)--(11) in \cite{Sim18}, p.~99; the limiting density $1/2$ requires the subsequent blow-up argument.
\end{theorem}
\begin{proof}
Use Proposition \ref{prop:reduced-boundary-perimeter-growth} and Proposition \ref{prop:phase-radial-sobolev} to obtain
$V^{(n-1)/n}\leq\mu(B_r)+V'\leq3V'$
for almost every small radius. Lemma \ref{lem:perimeter-support-two-phases} supplies positivity at every radius, and Lemma \ref{lem:radial-phase-volume} supplies absolute continuity and $V(0)=0$. Apply Lemma \ref{lem:phase-volume-ode} with $C=3$. The argument for $W$ uses the other flux bound and gives the same constant. The intersection of the full-measure sets of good radii suffices for the differential inequalities; their integrated conclusions hold for every small radius.
\end{proof}
